% Reverse-drafting note: contributions -> gaps -> challenges -> background.
\section{Introduction}
\label{sec:introduction}

Large panoramic-image collections make it possible to study indoor geometry at building scale, but they also move annotation quality from a drawing problem to a production-policy problem. A valid room layout must remain coherent across the equirectangular seam, preserve ceiling--wall and floor--wall ordering, and describe a plausible closed structure. Modern estimators such as HorizonNet and HoHoNet can provide useful initial layouts \cite{sun2019horizonnet,sun2021hohonet}, yet a usable reference still depends on human inspection and correction. The scale of resources such as ZInD illustrates both the scientific value and the substantial labour cost of panoramic layout annotation \cite{cruz2021zind}.

Model assistance does not make workers interchangeable. One worker may recover accurate boundaries but occasionally submit a structurally invalid polygon; another may agree consistently with peers while sharing a systematic deviation from the operational reference. A third may be globally reliable but less effective on a particular failure family. Routing every task by one overall score ignores this structure. Routing by a highly granular profile creates the opposite risk: an apparently favourable worker--task match may be an artefact of sparse calibration data. The operational question is therefore not simply whether workers differ, but \emph{when the available calibration evidence is strong enough for those differences to change routing}.

Three gaps motivate this study. First, research on panoramic layout estimation and efficient model-assisted annotation has largely optimized prediction, aggregation, stopping, or annotation effort \cite{sun2019horizonnet,sun2021hohonet,jiang2022lgtnet,branson2017lean,vanhorn2018leanmulticlass,liao2021goodpractices}. These works establish that machine predictions and worker models can reduce annotation cost, but they do not provide an evidence contract for routing geometric correction tasks using distinct reference, peer, and structural outcomes. The gap is not the absence of human--machine annotation systems; it is the absence of an auditable connection between multidimensional performance evidence and a deployed panoramic-layout routing decision.

Second, classic and recent crowdsourcing research shows that a single constant worker-accuracy parameter is often inadequate. Dawid and Skene separate latent truth from observer error \cite{dawid1979observer}; annotator-expertise models allow performance to vary over the input space \cite{yan2010expertise}; probabilistic models distinguish difficulty from subjectivity \cite{jin2018subjectivity}; and direct empirical tests reject several convenient constant-accuracy assumptions \cite{burrell2023conventional}. What remains unresolved for the present setting is how to estimate a small, interpretable profile when gold-like reference tasks are sparse, peer disagreement may be legitimate, and structural invalidity has several possible causes. We address this as a measurement problem, not as a claim that task-dependent worker ability is itself new.

Third, adaptive assignment and task recommendation already exploit worker heterogeneity \cite{ho2013adaptive,kang2018recommendation}, and recent online assignment work explicitly models reusable agents, rejection, and finite operational resources \cite{sumita2022reusable}. However, a prospective deployment needs a rule for unsupported worker--task combinations and an evaluation that does not favour one policy operationally. Adjacent work on safe policy improvement and support-deficient off-policy learning formalizes related concerns about uncertainty, coverage, and fallback \cite{laroche2019spibb,sachdeva2020support}. Our use of a global fallback is therefore not presented as the invention of safe fallback. The gap is narrower: no prior work cited here couples independent panoramic-layout calibration, an explicit support boundary, conditional routing, and a symmetric prospective policy test.

The setting creates three practical challenges. Measurement must distinguish accuracy against an operational reference from compatibility with stable peer solutions and from invalid output attributable to the worker. Calibration must transfer across stages without using the later policy outcomes that it is intended to predict. Finally, the routing policy must degrade predictably when conditional evidence is missing, ambiguous, version-inconsistent, or unstable. Treating these challenges jointly changes the unit of contribution from a ranking formula to an auditable chain of evidence.

We ask three research questions:

\begin{enumerate}
  \item \textbf{RQ1: Characterization.} Do reference-based accuracy, stable peer compatibility, and worker-attributable structural invalidity provide non-redundant and sufficiently stable descriptions of worker performance across independent calibration stages?
  \item \textbf{RQ2: Supported heterogeneity.} Which predeclared risk or failure families have enough calibration support to justify a conditional worker-performance component, and how reliably does the support rule suppress unsupported adjustments?
  \item \textbf{RQ3: Policy value.} Under symmetric prospective operations, does calibration-supported conditional routing preserve safety and improve delivery-adjusted quality or resource use relative to calibrated global ranking?
\end{enumerate}

The paper makes three evidence-bounded contributions. First, we operationalize a \emph{three-component worker performance characterization}: reference-based accuracy, stable peer compatibility, and worker-attributable structural invalidity. This measurement separates outcomes that have different meanings rather than collapsing them into an unrestricted latent ``ability.'' Second, we adapt independent cross-stage calibration, partial pooling, and support diagnostics to determine where a conditional component is estimable. This produces a calibration-support boundary and auditable distance from that boundary. Third, we propose \emph{calibration-supported conditional routing}: a policy that begins with a calibrated global ranking, activates only supported routable failure modes, and otherwise falls back to the global policy. Its main test is a prospective, 50:50 V1 policy experiment with symmetric constraints and ordered safety and quality endpoints; a separate T1 experiment supplies supporting evidence about model-assistance heterogeneity but is not treated as direct validation of V1.

The resulting claim is deliberately conditional. The framework does not imply that more personalized routing is always better, that peer agreement is truth, or that active time measures worker competence. It specifies what must be measured, frozen, and supported before worker--task heterogeneity is allowed to affect deployment. Section~\ref{sec:related} positions this claim; Sections~\ref{sec:setup}--\ref{sec:evaluation} define the design and analysis; Sections~\ref{sec:results} and \ref{sec:discussion} provide the marked result scaffold and its interpretation; and Section~\ref{sec:conclusion} states the intended evidence boundary.
